% Unit 047: English translation of chapter3.tex, lines 3201--3425.
% Source: Methods of Algebra, Volume 2, commit
% 9a5803ff77dd3257484cb177f851a73770a59dd3 (CC BY 4.0).
% stable-unit-id: o014.aljabr2.chapter3.exercises
% Coverage: complete Chapter 3 exercises (26 exercises and 17 hints).

% segment-id: o014.li2.chapter3.exercises.x001
\hypertarget{o014.aljabr2.chapter3.exercises}{ }
\begin{Exercises}
% segment-id: o014.li2.chapter3.exercises.ex001
	\item Let $\Bbbk$ be a commutative ring. Thus
	$\cate{C}(\Bbbk\dcate{Mod})$ and $\cate{K}(\Bbbk\dcate{Mod})$ are
	naturally $\Bbbk\dcate{Mod}$-categories in the sense of
	Definition~\ref{def:cat-k-linear}. Prove that, for every $n\in\Z$ and
	every object $X$ of $\cate{C}(\Bbbk\dcate{Mod})$, there are canonical
	isomorphisms of $\Bbbk$-modules
% segment-id: o014.li2.chapter3.exercises.q001
	\[
		\Hm^n(X) \simeq
		\Hom_{\cate{K}(\Bbbk\dcate{Mod})}(\Bbbk,X[n])
		\simeq
		\Hom_{\cate{K}(\Bbbk\dcate{Mod})}(\Bbbk[-n],X).
	\]
	On the right, $\Bbbk$ is regarded as a complex concentrated in degree
	$0$.

% segment-id: o014.li2.chapter3.exercises.ex002
	\item This exercise relates homotopies between morphisms to isomorphisms
	between their mapping cones. Let $\mathcal{A}$ be an additive category and
	take any pair of morphisms
	$f,g\in\Hom_{\cate{C}(\mathcal{A})}(X,Y)$. Construct a bijection
% segment-id: o014.li2.chapter3.exercises.q002
	\[
	\left\{
	\begin{array}{c}
		s\in\Hom^{-1}(X,Y): \\
		d^{-1}_{\Hom^\bullet(X,Y)}(s)=f-g
	\end{array}
	\right\}
	\xrightarrow{1:1}
	\left\{
	\begin{array}{c}
		h:\Cone(f)\to\Cone(g),\\
		\text{a morphism of complexes making}\\
		\text{the following diagram commute}
	\end{array}
	\right\},
	\]
% segment-id: o014.li2.chapter3.exercises.g001
	\begin{equation*}
	\begin{tikzcd}[row sep=small]
		& \Cone(f) \arrow[dr, "\beta(f)"] \arrow[dd, "h"] & \\
		Y \arrow[ru, "\alpha(f)"] \arrow[rd, "\alpha(g)"'] & & X[1] \\
		& \Cone(g) \arrow[ru, "\beta(g)"'] &
	\end{tikzcd}
	\end{equation*}
	and prove that every such morphism $h$ is an isomorphism of complexes. In
	particular, $f$ and $g$ are homotopic if and only if such an $h$ exists.

% segment-id: o014.li2.chapter3.exercises.hn001
	\begin{hint}
		Commutativity of the diagram in degree $n$ is equivalent to the matrix
		equation
% segment-id: o014.li2.chapter3.exercises.q003
		\[
			h^n=
			\begin{pmatrix}
				\identity_{X^{n+1}}&0\\
				s^{n+1}&\identity_{Y^n}
			\end{pmatrix},
			\quad\text{where}\quad
			s^{n+1}:X^{n+1}\to Y^n.
		\]
	\end{hint}

% segment-id: o014.li2.chapter3.exercises.ex003
	\item (Split exactness) Let $\mathcal{A}$ be an Abelian category and $X$
	a complex over it. Prove that the following statements are equivalent.
% segment-id: o014.li2.chapter3.exercises.l001
	\begin{enumerate}[(i)]
% segment-id: o014.li2.chapter3.exercises.i001
		\item $X$ is equal to $0$ in $\cate{K}(\mathcal{A})$;
% segment-id: o014.li2.chapter3.exercises.i002
		\item $\identity_X$ is null-homotopic;
% segment-id: o014.li2.chapter3.exercises.i003
		\item $X$ is acyclic and there exists
		$(s^n)_n\in\Hom^{-1}(X,X)$ such that
		$d^ns^{n+1}d^n=d^n$ for every $n$;
% segment-id: o014.li2.chapter3.exercises.i004
		\item there is a family of objects $(Y^n)_n$ in $\mathcal{A}$ and an
		isomorphism of complexes
		$\Phi:X\rightiso(Y^n\oplus Y^{n+1},\overline d^{\,n})_n$, where
		$\overline d$ is represented by the matrix
		$\bigl(\begin{smallmatrix}0&1\\0&0\end{smallmatrix}\bigr)$.
	\end{enumerate}

% segment-id: o014.li2.chapter3.exercises.p001
	A complex satisfying any of these conditions is also called
	\emph{split exact}.\footnote{More generally, let $X$ be a complex over an
	additive category $\mathcal{A}$. If there are two families of objects
	$(Y^n)_n$ and $(H^n)_n$ such that
	$X^n\simeq Y^n\oplus H^n\oplus Y^{n+1}$ and $d$ is identified with
	$\bigl(\begin{smallmatrix}0&0&1\\0&0&0\\0&0&0\end{smallmatrix}\bigr)$,
	then $X$ is called \emph{split}.}
	For
	$X=[\cdots\,0\to A\to B\to C\to0\,\cdots]$, prove that this condition
	is equivalent to $X$ being a split short exact sequence.
	\index{complex!split exact}

% segment-id: o014.li2.chapter3.exercises.hn002
	\begin{hint}
		It suffices to prove (iii) $\implies$ (iv). The condition implies that
		$d^{n-1}s^n\in\End_{\mathcal{A}}(X^n)$ is idempotent. This gives a
		decomposition $X^n=Y^n\oplus Z^n$ with
		$Y^n=\Image(d^{n-1}s^n)$. Prove that
		$Y^n=\Image(d^{n-1})$ and that the restriction of $d^n$ gives an
		isomorphism $Z^n\rightiso\Image(d^n)$; then take $\Phi$ corresponding
		to $\bigl(\begin{smallmatrix}ds\\d\end{smallmatrix}\bigr)$.
	\end{hint}

% segment-id: o014.li2.chapter3.exercises.ex004
	\item Let $\mathcal{A}$ be an Abelian category. Prove that $X$ is an
	injective object of $\cate{C}(\mathcal{A})$ if and only if every term of
	$X$ is injective in $\mathcal{A}$ and $X$ is equal to $0$ in
	$\cate{K}(\mathcal{A})$. State the version for projective objects.

% segment-id: o014.li2.chapter3.exercises.hn003
	\begin{hint}
		For the only-if direction, observe that the short exact sequence in
		$\cate{C}(\mathcal{A})$
		$0\to X\to\Cone(\identity_X)\to X[1]\to0$ splits.
		Proposition~\ref{prop:cone-homotopy} then shows that $X$ is acyclic.
		Choose a section $X[1]\to\Cone(\identity_X)$ and write it as
		$\bigl(\begin{smallmatrix}1\\-\theta\end{smallmatrix}\bigr)$, where
		$\theta^n:X^{n+1}\to X^n$. Verify that
		$d_X^{n-1}\theta^{n-1}+\theta^nd_X^n=\identity_{X^n}$ to prove that
		$\identity_X$ is null-homotopic.

		Fix $n$. By the preceding exercise, decompose
		$X^n=Y^n\oplus Y^{n+1}$; it remains to prove that $Y^n$ is injective.
		For any morphism $f:A\to Y^n$ and monomorphism
		$g:A\hookrightarrow B$ in $\mathcal{A}$, consider the commutative
		diagram with exact row in $\cate{C}(\mathcal{A})$
% segment-id: o014.li2.chapter3.exercises.g002
		\[
		\begin{tikzcd}[column sep=large]
			0 \arrow[r] &
			{A[-n]} \arrow[r, "{g[-n]}"] \arrow[d, "\Psi"'] &
			{B[-n]} \arrow[dashed, ld, "\exists"] \\
			& X &
		\end{tikzcd}
		\]
		where $\Psi^n=\bigl(\begin{smallmatrix}f\\0\end{smallmatrix}\bigr)$.
		Use this diagram to verify the injectivity criterion for $Y^n$.

		For the if direction, again write $X^n=Y^n\oplus Y^{n+1}$ and
		describe explicitly all possible morphisms $f:Z\to X$.
	\end{hint}

% segment-id: o014.li2.chapter3.exercises.ex005
	\item Let $\Bbbk$ be a commutative ring and let $R$ be a unital
	$\Bbbk$-algebra, in accordance with the convention of this book. Write
	down the canonical isomorphisms
	$\HHm_0(M_n(R))\simeq\HHm_0(R)$ and
	$\HHm^0(M_n(R))\simeq\HHm^0(R)$, where $M_n(R)$ is the
	$\Bbbk$-algebra of $n\times n$ matrices over $R$. In fact, analogous
	isomorphisms hold in every degree $m$, for both $\HHm_m$ and $\HHm^m$;
	see \cite[Theorem~1.2.4]{Lo98}.

% segment-id: o014.li2.chapter3.exercises.ex006
	\item In the setting of \S\ref{sec:HH}, define the sub-bimodule
	$\mathrm{Triv}_nR$ of $\mathsf{B}_nR$, for every
	$n\in\Z_{\geq0}$, by
% segment-id: o014.li2.chapter3.exercises.q004
	\begin{align*}
		\mathrm{Triv}_nR
		&:=\sum_{h=1}^n
		R\otimes
		(\cdots\otimes
		\underbracket{\;\Bbbk\;}_{\text{$h$-th factor}}
		\otimes\cdots)\otimes R,
		\qquad \mathrm{Triv}_0R=0.
	\end{align*}
% segment-id: o014.li2.chapter3.exercises.l002
	\begin{enumerate}[(i)]
% segment-id: o014.li2.chapter3.exercises.i005
		\item Verify that this defines a subcomplex $\mathrm{Triv}R$ of
		$\mathsf{B}R$, and that $\mathrm{Triv}R$ is exact.
% segment-id: o014.li2.chapter3.exercises.hn004
		\begin{hint}
			Restrict $(h^n)_n$ from the proof of
			Lemma~\ref{prop:HH-bar-exactness} to $\mathrm{Triv}R$.
		\end{hint}
% segment-id: o014.li2.chapter3.exercises.i006
		\item Write $\overline R:=R/\Bbbk$ as a $\Bbbk$-module and define
		the \emph{reduced bar complex} of $R$ to be the quotient
% segment-id: o014.li2.chapter3.exercises.q005
		\[
			\overline{\mathsf{B}}R
			:=\mathsf{B}R/\mathrm{Triv}R
			=\left(R\otimes\overline R^{\otimes n}\otimes R,b_n\right)_{n\geq0}.
		\]
		Following the pattern of \eqref{eqn:HH-cplx}, simplify
		$M\dotimes{R^e}\overline{\mathsf{B}}R$ and
		$\Hom_{R^e}(\overline{\mathsf{B}}R,M)$.
% segment-id: o014.li2.chapter3.exercises.i007
		\item Give canonical isomorphisms
		$\Hm_n(M\dotimes{R^e}\overline{\mathsf{B}}R)\simeq\HHm_n(M)$ and
		$\Hm^n(\Hom_{R^e}(\overline{\mathsf{B}}R,M))\simeq\HHm^n(M)$.
% segment-id: o014.li2.chapter3.exercises.hn005
		\begin{hint}
			In general, this follows from the theorem on normalized complexes in
			the Dold--Kan correspondence,
			\sourcecrossref{prop:Dold-Kan-N-C}{in \S~8.5}, provided that
			$\mathsf{B}R$ is understood from the viewpoint of
			\sourcecrossref{eg:HH-comonad}{the Example in \S~8.7}.
		\end{hint}
	\end{enumerate}

% segment-id: o014.li2.chapter3.exercises.ex007
	\item Let $\Bbbk$ be a commutative ring, and let $\widetilde R$ and $R$
	be unital $\Bbbk$-algebras, in accordance with the convention of this
	book. Suppose there is a split short exact sequence of $\Bbbk$-modules
% segment-id: o014.li2.chapter3.exercises.g003
	\[
	\begin{tikzcd}
		0 \arrow[r] & M \arrow[r] & \widetilde R
		\arrow[shift left, r, "p"] &
		R \arrow[r] \arrow[shift left, l, "s"] & 0,
	\end{tikzcd}
	\]
	where $M$ is a two-sided ideal of $\widetilde R$, $ps=\identity_R$, and
	$M^2=\{0\}$. Use $s$ to identify $\widetilde R$ with $R\oplus M$.
% segment-id: o014.li2.chapter3.exercises.l003
	\begin{enumerate}[(i)]
% segment-id: o014.li2.chapter3.exercises.i008
		\item For $m\in M$ and $r\in R$, choose any
		$\widetilde r\in p^{-1}(r)$. Prove that $\widetilde rm$ and
		$m\widetilde r$ depend only on $m$ and $r$, and that these operations
		make $M$ an $(R,R)$-bimodule.
% segment-id: o014.li2.chapter3.exercises.i009
		\item Prove that there is a bilinear map $f:R^2\to M$ such that
		multiplication in $\widetilde R$ is given by
% segment-id: o014.li2.chapter3.exercises.q006
		\[
			(r_1,m_1)(r_2,m_2)
			=(r_1r_2,r_1m_2+m_1r_2+f(r_1,r_2)),
			\quad r_1,r_2\in R,\quad m_1,m_2\in M.
		\]
% segment-id: o014.li2.chapter3.exercises.i010
		\item Such data $(R,\widetilde R,M)$ are called a
		\emph{square-zero extension} of $R$ by $M$. Prove that, up to the
		appropriate notion of isomorphism, square-zero extensions are
		classified by $\HHm^2(M)$, with the zero element corresponding to the
		trivial extension $R\times M$.
% segment-id: o014.li2.chapter3.exercises.hn006
		\begin{hint}
			Compute $\HHm^2(M)$ using the reduced bar complex from the preceding
			exercise.
		\end{hint}
	\end{enumerate}

% segment-id: o014.li2.chapter3.exercises.ex008
	\item Complete the proof of Lemma~\ref{prop:CC-cplx}.

% segment-id: o014.li2.chapter3.exercises.ex009
	\item In the setting of Definition~\ref{def:HC}, verify for the special
	case $R=\Bbbk$ that:
% segment-id: o014.li2.chapter3.exercises.l004
	\begin{enumerate}[(i)]
% segment-id: o014.li2.chapter3.exercises.i011
		\item if $n$ is odd, then
		$\mathrm{HP}_n(\Bbbk)=\mathrm{HC}_n(\Bbbk)=0$;
% segment-id: o014.li2.chapter3.exercises.i012
		\item if $n$ is even (or nonnegative even), then
		$\mathrm{HP}_n(\Bbbk)$ (or $\mathrm{HC}_n(\Bbbk)$) is isomorphic to
		$\Bbbk$.
	\end{enumerate}

% segment-id: o014.li2.chapter3.exercises.ex010
	\item In the setting of Definition~\ref{def:HC} and
	Theorem~\ref{prop:Connes-exact}, take the polynomial ring $R=\Bbbk[t]$.
% segment-id: o014.li2.chapter3.exercises.l005
	\begin{enumerate}[(i)]
% segment-id: o014.li2.chapter3.exercises.i013
		\item Prove that $\Omega_{R\mid\Bbbk}$ in Example~\ref{eg:HH1} is a
		free $R$-module of rank $1$, with basis $\dd t$.
% segment-id: o014.li2.chapter3.exercises.i014
		\item Describe the connecting morphism
		$B:\mathrm{HC}_0(R)\to\mathrm{HH}_1(R)$; note that both sides are
		isomorphic to $R$.
% segment-id: o014.li2.chapter3.exercises.i015
		\item Describe all the $\mathrm{HC}_n(R)$ as explicitly as possible.
		For $\Bbbk=\Z$, prove that
		$\mathrm{HC}_1(\Z[t])=\bigoplus_{n\geq2}\Z/n\Z$.
	\end{enumerate}
% segment-id: o014.li2.chapter3.exercises.hn007
	\begin{hint}
		Use the long exact sequence of Theorem~\ref{prop:Connes-exact} and the
		calculation of $\HHm_n(\Bbbk[t])$ in high degrees in
		Example~\ref{eg:HH-Tor}.
	\end{hint}

% segment-id: o014.li2.chapter3.exercises.ex011
	\item Let $\mathcal{A}$ be an Abelian category and let
	$0\to X\to I^0\to I^1\to\cdots$ be an exact sequence in
	$\cate{C}(\mathcal{A})$. Regard $I=[I^0\to I^1\to\cdots]$ as a
	double complex, with $I^{p,q}:=(I^p)^q$. Assuming that
	$I\in\Obj(\cate{C}^2_f(\mathcal{A}))$, prove that the evident morphism
	$X\to\tot(I)$ is a quasi-isomorphism.

% segment-id: o014.li2.chapter3.exercises.ex012
	\item Let $X\in\Obj(\cate{C}^2_f(\mathcal{A}))$, where
	$\mathcal{A}$ is an Abelian category. Prove that, if the $i$-th column
	$\left(X^{i,\bullet},\dvert\right)$ and the $j$-th row
	$\left(X^{\bullet,j},\dhori\right)$ are exact for every
	$i,j\in\Z\smallsetminus\{0\}$, then for every $n\in\Z$ there is a
	canonical isomorphism
	$\Hm^n(X^{0,\bullet},\dvert)\simeq
	\Hm^n(X^{\bullet,0},\dhori)$.

% segment-id: o014.li2.chapter3.exercises.hn008
	\begin{hint}
		Define the brutal truncation functor
		$\sigma_{\mathrm{I}}^{\leq n}$ (respectively
		$\sigma_{\mathrm{I}}^{\geq n}$) from
		$\cate{C}^2(\mathcal{A})$ to itself by replacing the terms
		$(X^{i,j})_{(i,j)\in\Z^2}$ for $i>n$ (respectively $i<n$) by $0$ and
		leaving all other terms unchanged. Use
		Theorem~\ref{prop:double-cplx-tot} to prove that both morphisms in
% segment-id: o014.li2.chapter3.exercises.q007
		\[
			\text{$0$-th column}
			=\sigma_{\mathrm{I}}^{\leq0}\sigma_{\mathrm{I}}^{\geq0}(X)
			\twoheadleftarrow
			\sigma_{\mathrm{I}}^{\geq0}X
			\hookrightarrow X
		\]
		induce quasi-isomorphisms on total complexes. Do the same after
		interchanging the roles of rows and columns.
	\end{hint}

% segment-id: o014.li2.chapter3.exercises.ex013
	\item Use Lemma~\ref{prop:ext-alpha-beta} to prove the existence part for
	$\beta$ in Theorem~\ref{prop:resolution-homotopy}.

% segment-id: o014.li2.chapter3.exercises.hn009
	\begin{hint}
		It suffices to treat the first case. The mapping-cylinder lemma,
		\ref{prop:cone-beta}, factors $X\xrightarrow{\alpha}Y$ as
		$X\xrightarrow{\alpha'}\Cyl(\alpha)\xrightarrow{\psi}Y$, where
		$\alpha'$ is a monomorphism and $\psi$ is invertible in
		$\cate{K}(\mathcal{A})$.
	\end{hint}

% segment-id: o014.li2.chapter3.exercises.ex014
	\item With the notation of the preceding exercise, prove the uniqueness
	part for $\beta$ in Theorem~\ref{prop:resolution-homotopy}.

% segment-id: o014.li2.chapter3.exercises.hn010
	\begin{hint}
		It suffices to treat the first case. Use the mapping cylinder again to
		reduce to the case in which every
		$\alpha^n:X^n\hookrightarrow Y^n$ is the inclusion of a direct
		summand. Suppose the morphisms $\beta_i:Y\to I$ in
		$\cate{C}(\mathcal{A})$, for $i=1,2$, satisfy
		$\gamma-\beta_i\alpha=d^{-1}_{\Hom^\bullet(X,I)}h_i$, where
		$h_i\in\Hom^{-1}(X,I)$. Let $C:=\Coker(\alpha)$ be the acyclic complex
		and define $h^n:Y^n\simeq X^n\oplus C^n\to I^{n-1}$ so that its
		restriction to $X^n$ (respectively $C^n$) is $h_1^n-h_2^n$
		(respectively $0$). Verify that
% segment-id: o014.li2.chapter3.exercises.q008
		\[
			\beta_1-\beta_2-d^{-1}_{\Hom^\bullet(Y,I)}h:Y\to I
		\]
		becomes $0$ after composition with $\alpha$. Thus, up to homotopy,
		$\beta_1-\beta_2$ factors as $Y\twoheadrightarrow C\to I$; apply
		Lemma~\ref{prop:resolution-homotopy-aux} to the second morphism.
	\end{hint}

% segment-id: o014.li2.chapter3.exercises.ex015
	\item Let $\mathcal{A}$ be an Abelian category. Prove that
	$(\Hm^n)_{n\geq0}$ is a universal cohomological $\delta$-functor from the
	Abelian category $\cate{C}^{\geq0}(\mathcal{A})$ (see
	Definition~\ref{def:cplx-cat-variant}) to $\mathcal{A}$.
% segment-id: o014.li2.chapter3.exercises.hn011
	\begin{hint}
		To show that $\Hm^n$ is effaceable when $n>0$, take
		$\alpha(\identity_X):X\hookrightarrow\Cone(\identity_X)$. This
		morphism factors through the brutal truncation
		$\sigma^{\geq0}\Cone(\identity_X)$. Use
		Proposition~\ref{prop:cone-homotopy} (i).
	\end{hint}

% segment-id: o014.li2.chapter3.exercises.ex016
	\item (Complex version of the Milnor exact sequence) Let $R$ be any ring
	and let $(X_k,f_k)_{k\geq0}$ be an object of
	$\cate{InvSys}(\cate{C}(R\dcate{Mod}))$. Suppose
	$(X_k^n,f_k^n)_{k\geq0}$ satisfies the Mittag--Leffler condition for every
	$n\in\Z$. The notation $f_k$ will henceforth be suppressed. Follow these
	steps to construct a short exact sequence
% segment-id: o014.li2.chapter3.exercises.q009
	\[
		0\to\lim\nolimits^1_k\Hm^{n-1}(X_k)
		\to\Hm^n(\varprojlim_kX_k)
		\to\varprojlim_k\Hm^n(X_k)\to0.
	\]
% segment-id: o014.li2.chapter3.exercises.l006
	\begin{enumerate}[(i)]
% segment-id: o014.li2.chapter3.exercises.i016
		\item As in Theorem~\ref{prop:CE-resolution}, define
		$B_k^n\subset Z_k^n\subset X_k^n$ and $H_k^n$, and make
		$B_k\subset Z_k\subset X_k$ and $H_k$ complexes with
		$d_{B_k}=d_{Z_k}=d_{H_k}=0$. Prove that
		$(B_k^n)_{k\geq0}$ satisfies the Mittag--Leffler condition for every
		$n$.
% segment-id: o014.li2.chapter3.exercises.i017
		\item Prove that
		$0\to\varprojlim_kZ_k\to\varprojlim_kX_k
		\xrightarrow{d}(\varprojlim_kX_k)[1]$ is exact.
% segment-id: o014.li2.chapter3.exercises.i018
		\item Define
		$B^n:=\Image(d^{n-1}:\varprojlim_kX_k^{n-1}\to
		\varprojlim_kX_k^n)$\footnote{Translator's note (O014-C067): the
		source omits the subscript $k$ from the limit in the codomain, even
		though the domain, the system under discussion, and the subsequent
		exact sequence all use $\varprojlim_k$. The target restores that
		subscript.}, and make $B$ a complex with $d_B=0$. Prove that there is a
		short exact sequence
		$0\to B[1]\to\varprojlim_kB_k[1]\to\lim\nolimits^1_kZ_k\to0$.
% segment-id: o014.li2.chapter3.exercises.hn012
		\begin{hint}
			Use the short exact sequence
			$0\to Z_k\to X_k\xrightarrow{d}B_k[1]\to0$ and the
			Mittag--Leffler condition on $(X_k^n)_{k\geq0}$.
		\end{hint}
% segment-id: o014.li2.chapter3.exercises.i019
		\item Deduce the isomorphism
		$\lim\nolimits^1_kZ_k\rightiso\lim\nolimits^1_kH_k$ and the short
		exact sequence
		$0\to\varprojlim_kB_k\to\varprojlim_kZ_k
		\to\varprojlim_kH_k\to0$.
% segment-id: o014.li2.chapter3.exercises.hn013
		\begin{hint}
			Use the short exact sequence
			$0\to B_k\to Z_k\to H_k\to0$ and the Mittag--Leffler condition on
			$(B_k^n)_{k\geq0}$.
		\end{hint}
% segment-id: o014.li2.chapter3.exercises.i020
		\item Take the appropriate inclusions
		$B\subset\varprojlim_kB_k\subset\varprojlim_kZ_k$ and consider the
		corresponding quotients to obtain the required short exact sequence.
	\end{enumerate}

% segment-id: o014.li2.chapter3.exercises.ex017
	\item Let $\Bbbk$ be a field. Define objects $A$ and $B$ of
	$\cate{InvSys}(\Bbbk\dcate{Mod})$ by
	$A_k:=X^k\Bbbk\llbracket X\rrbracket$ and
	$B_k:=X^k\Bbbk[X]$, taking the transition morphisms to be the inclusions.
	Neither satisfies the Mittag--Leffler condition. Prove that
	$\lim\nolimits^1A=0$, whereas
	$\lim\nolimits^1B\simeq
	\Bbbk\llbracket X\rrbracket/\Bbbk[X]$.
% segment-id: o014.li2.chapter3.exercises.hn014
	\begin{hint}
		Embed these systems of ideals into systems of rings, and use the
		dimension-shifting technique of Proposition~\ref{prop:dimension-shifting}
		to compute $\lim\nolimits^1$.
	\end{hint}

% segment-id: o014.li2.chapter3.exercises.ex018
	\item Let $D$ be a division ring. Prove that $\Ext_D^n$ and $\Tor_n^D$
	vanish for $n>0$.

% segment-id: o014.li2.chapter3.exercises.ex019
	\item For every ring $R$, formulate precisely and prove the canonical
	isomorphism $\Tor_n^R(X,Y)\simeq\Tor_n^{R^{\opp}}(Y,X)$.

% segment-id: o014.li2.chapter3.exercises.ex020
	\item Let the integral domain $R$ be a principal ideal domain. For all
	finitely generated $R$-modules $M,N$ and every $n\in\Z_{\geq0}$,
	describe $\Ext_R^n(M,N)$ and $\Tor_n^R(M,N)$ explicitly.

% segment-id: o014.li2.chapter3.exercises.ex021
	\item Let $M$ be a $\Z$-module. Prove that $\Tor_1^\Z(\Q/\Z,M)$ is the
	subgroup of $M$ consisting of all its torsion elements.

% segment-id: o014.li2.chapter3.exercises.ex022
	\item (Universal coefficient theorem for cohomology) Let $R$ be a ring
	and let $C=(C_\bullet,d^C_\bullet)$ be a chain complex of left
	$R$-modules. For any left $R$-module $M$, applying
	$\Hom_R(\cdot,M)$ termwise gives a cochain complex
	$C^\bullet:=\Hom(C_\bullet,M)$; denote its degree-$n$ cohomology by
	$\Hm^n(C,M)$. Suppose that
	$B_n:=\Image(d_{n+1})$ and $Z_n:=\Ker(d_n)$ are projective modules for
	every $n\in\Z$, and define $H_n:=Z_n/B_n=\Hm_n(C)$. Prove that, for
	every $n\in\Z$, there is a canonical short exact sequence
% segment-id: o014.li2.chapter3.exercises.q010
	\[
		0\to\Ext_R^1(\Hm_{n-1}(C),M)
		\to\Hm^n(C,M)
		\to\Hom_R(\Hm_n(C),M)\to0.
	\]
	Prove also that, if $R$ is a principal ideal domain and every $C_n$ is
	free, then these hypotheses always hold.
	\index{universal coefficient theorem}

% segment-id: o014.li2.chapter3.exercises.hn015
	\begin{hint}
		Define the functor $D(\cdot):=\Hom_R(\cdot,M)$ from
		$(R\dcate{Mod})^{\opp}$ to $\cate{Ab}$. Use the short exact sequence
		$0\to H_n\to C_n/B_n\to B_{n-1}\to0$ and the hypotheses to obtain a
		commutative diagram with exact rows
% segment-id: o014.li2.chapter3.exercises.g004
		\[
		\begin{tikzcd}
			0 \arrow[r] &
			D(Z_{n-1}) \arrow[r, "\identity"] \arrow[d] &
			D(Z_{n-1}) \arrow[r] \arrow[d, "{D(d^C_n)}"] &
			0 \arrow[r] \arrow[d] & 0 \\
			0 \arrow[r] &
			D(B_{n-1}) \arrow[r, "{D(d^C_n)}"' inner sep=0.5em] &
			D(C_n/B_n) \arrow[r] &
			D(H_n) \arrow[r] & 0 .
		\end{tikzcd}
		\]
		Prove that $d_{D(C)}^{n-1}$ factors as
		$D(C_{n-1})\twoheadrightarrow D(Z_{n-1})
		\xrightarrow{D(d^C_n)}D(C_n/B_n)\hookrightarrow D(C_n)$ and that
		$D(C_n/B_n)\simeq\Ker(d_C^n)$. Thus the cokernel of the middle
		morphism in the diagram is $\Hm^n(C,M)$. On the other hand, the
		sequence $0\to B_{n-1}\to Z_{n-1}\to H_{n-1}\to0$ shows that the
		cokernel of the left morphism is
		$\Ext_R^1(\Hm_{n-1}(C),M)$. Apply the snake lemma,
		Theorem~\ref{prop:snake-lemma}.
	\end{hint}

% segment-id: o014.li2.chapter3.exercises.ex023
	\item In the setting of \S\ref{sec:HH}, take
	$R=\Bbbk[t]/(t^n)$, where $t$ is a polynomial variable and
	$n\in\Z_{\geq1}$, so that
	$R^e=\Bbbk[x,y]/(x^n,y^n)$.\footnote{Translator's note (O014-C068): the
	source gives the coefficient ring as $R[x,y]$. However, for
	$R=\Bbbk[t]/(t^n)$, the ring
	$R^e=R\otimes_{\Bbbk}R^{\opp}$ is canonically
	$\Bbbk[x,y]/(x^n,y^n)$. The target restores the required coefficient ring
	$\Bbbk$.}
% segment-id: o014.li2.chapter3.exercises.l007
	\begin{enumerate}
% segment-id: o014.li2.chapter3.exercises.i021
		\item Prove that the following chain complex is exact:
% segment-id: o014.li2.chapter3.exercises.q011
		\[
			\cdots\xrightarrow{v}R^e\xrightarrow{u}R^e
			\xrightarrow{v}R^e\xrightarrow{u}R^e\to R\to0,
		\]
		where $R^e\to R$ maps $f(x,y)$ to $f(t,t)$ and
% segment-id: o014.li2.chapter3.exercises.q012
		\[
			u:=x-y,\qquad
			v:=\sum_{\substack{a+b=n-1\\a,b\geq0}}x^ay^b.
		\]
		Deduce that, for every $R$-module $M$,
		$\HHm_{k+2}(M)\simeq\HHm_k(M)$ and
		$\HHm^{k+2}(M)\simeq\HHm^k(M)$.
% segment-id: o014.li2.chapter3.exercises.i022
		\item Let $\Bbbk$ be a field. Describe $\HHm_k(M)$ and $\HHm^k(M)$,
		distinguishing the cases $k\geq1$ and whether
		$\mathrm{char}(\Bbbk)$ is relatively prime to $n$.
	\end{enumerate}

% segment-id: o014.li2.chapter3.exercises.ex024
	\item For any ring $R$ and $n\in\Z_{\geq0}$, prove that $\Tor_n^R$
	commutes with small filtered inductive limits; that is,
% segment-id: o014.li2.chapter3.exercises.q013
	\[
		\Tor_n^R(\varinjlim_iX_i,Y)
		\simeq\varinjlim_i\Tor_n^R(X_i,Y),
		\qquad
		\Tor_n^R(X,\varinjlim_iY_i)
		\simeq\varinjlim_i\Tor_n^R(X,Y_i).
	\]
	Discuss whether there is an analogous statement for the $\Ext$ functors.

% segment-id: o014.li2.chapter3.exercises.ex025
	\item Prove that an object $I$ of an Abelian category $\mathcal{A}$ is
	injective if and only if, regarded as a complex, $I$ is K-injective. State
	the dual version.
% segment-id: o014.li2.chapter3.exercises.hn016
	\begin{hint}
		One direction follows from Example~\ref{eg:bdd-K-resolution}. For the
		converse, regard a short exact sequence
		$0\to A\to B\to C\to0$ in $\mathcal{A}$ as an acyclic complex, and
		construct a morphism from that complex to $I$ from the given morphism
		$A\to I$.
	\end{hint}

% segment-id: o014.li2.chapter3.exercises.ex026
	\item (N.\ Nitsure) Let $F:\mathcal{A}\to\mathcal{B}$ be a left exact
	functor between Abelian categories. Suppose $\mathcal{A}$ has enough
	injective objects and $X\in\Obj(\mathcal{A})$. Take an injective resolution
	$0\to X\to I^0\to I^1\to\cdots$ and an $F$-acyclic resolution
	$0\to X\to A^0\to A^1\to\cdots$. Write
% segment-id: o014.li2.chapter3.exercises.q014
	\[
		I=[\cdots\to0\to I^0\to I^1\to\cdots],
		\qquad
		A=[\cdots\to0\to A^0\to A^1\to\cdots].
	\]
% segment-id: o014.li2.chapter3.exercises.l008
	\begin{itemize}
% segment-id: o014.li2.chapter3.exercises.i023
		\item Lemma~\ref{prop:resolution-homotopy-classical} determines a
		unique morphism $\beta:A\to I$ in $\cate{K}(\mathcal{A})$ compatible
		with $A\leftarrow X\rightarrow I$. After applying $F$, this morphism
		gives\footnote{Translator's note (O014-C069): the source writes
		$\Hm^n(\beta):\Hm^n(A)\to\Hm^n(I)$ and, in the following item,
		$\Hm^n(A)$. Since $F$-acyclicity is being used to compute the derived
		functors of $F$, the morphism and complexes whose cohomology must be
		taken are $F(\beta)$, $F(A)$, and $F(I)$. The target restores the
		application of $F$.}
		\[
			c^n:=\Hm^n(F(\beta)):
			\Hm^n(F(A))\to\Hm^n(F(I))=\mathrm{R}^nF(X).
		\]
% segment-id: o014.li2.chapter3.exercises.i024
		\item On the other hand, Corollary~\ref{prop:acyclic-resolution} gives,
		by dimension shifting, an isomorphism
		\[
			d^n:\Hm^n(F(A))\rightiso\mathrm{R}^nF(X).
		\]
	\end{itemize}
	Prove that $d^n=(-1)^{n(n+1)/2}c^n$ for every $n\geq0$.
% segment-id: o014.li2.chapter3.exercises.hn017
	\begin{hint}
		This is the content of \cite{Ni09}.
	\end{hint}
\end{Exercises}
